\subsection{Contoh Perilaku Demokrasi Pancasila}

\subsubsection{Lingkungan Sekolah}
\begin{itemize}
\item Aktif dalam kegiatan diskusi kelas.
\item Menyampaikan aspirasi secara santun.
\item Menggunakan hak pilih dalam pemilihan Ketua OSIS.
\item Mengembangkan musyawarah mufakat dalam mengatasi permasalahan.
\item Berani mengakui kekalahan dalam pemilihan ketua kelas atau Ketua OSIS.
\end{itemize}

\subsubsection{Lingkungan Masyarakat}
\begin{itemize}
\item Antre secara tertib menunggu giliran.
\item Menghindari sikap main hakim sendiri dan tindakan kekerasan.
\item Menghormati para pemimpin, baik tokoh formal maupun adat.
\item Berani mengakui kebenaran atau kebaikan pendapat orang lain.
\item Melaksanakan hasil musyawarah warga dengan penuh tanggung jawab.
\end{itemize}

\subsubsection{Lingkungan Bangsa dan Negara}
\begin{itemize}
\item Mematuhi aturan hukum yang berlaku.
\item Menyampaikan aspirasi secara santun dan tertib.
\item Menggunakan hak pilih secara bertanggung jawab.
\item Menyukseskan program pemerintah yang sedang berkuasa.
\item Berani mengakui kekalahan dan memberikan ucapan selamat kepada pemenang pemilihan.
\end{itemize}

\subsubsection{Prinsip Kebebasan yang Bertanggung Jawab}
\begin{itemize}
\item Menggunakan kebebasan berpendapat dan berorganisasi secara bertanggung jawab.
\item Mendukung pelaksanaan Pemilihan Umum yang Langsung, Umum, Bebas, Rahasia, Jujur, dan Adil.
\item Menghormati orang lain dalam menyampaikan pendapat.
\item Menaati persyaratan dan hukum yang berlaku dalam melakukan kebebasan berorganisasi atau berserikat.
\end{itemize}

\subsection{Kemerdekaan Berpendapat dan Kebebasan Berekspresi}

\subsubsection{Pentingnya Kebebasan Berekspresi}
\begin{itemize}
\item Menjamin pemenuhan diri seseorang.
\item Mencari alternatif masukan dari berbagai pandangan yang berbeda.
\item Memungkinkan partisipasi warga negara dalam proses pengambilan keputusan politik.
\item Mendorong pencapaian stabilitas dan adaptasi dalam kehidupan masyarakat dan negara.
\end{itemize}

\subsubsection{Bentuk Penyampaian Pendapat dan Batasannya}
\begin{itemize}
\item Penyampaian secara lisan dapat dilakukan melalui pidato, dialog, diskusi, seminar, mimbar bebas, maupun demonstrasi.
\item Penyampaian secara tertulis dapat dilakukan melalui tulisan di media massa, media sosial, petisi, gambar, pamflet, poster, brosur, selebaran, dan spanduk.
\item Kebebasan berpendapat dibatasi oleh norma, hukum, serta larangan terhadap penyebaran berita bohong, fitnah, ujaran kebencian, dan aksi anarkis.
\end{itemize}

\subsubsection{Landasan Hukum dan Kewajiban Negara}
\begin{itemize}
\item Landasan hukum kebebasan berpendapat dan berorganisasi meliputi Undang-Undang Nomor 9 Tahun 1998 tentang Kemerdekaan Menyampaikan Pendapat di Muka Umum, Undang-Undang Nomor 40 Tahun 1999 tentang Pers, serta Undang-Undang Nomor 19 Tahun 2016 tentang Informasi dan Transaksi Elektronik.
\item Peraturan terkait kebebasan berserikat dan berkumpul diatur dalam Undang-Undang MD3, Undang-Undang Ormas, Undang-Undang Pemilu, dan Undang-Undang Partai Politik.
\item Negara memiliki tiga kewajiban utama terhadap hak asasi manusia, yaitu kewajiban menghormati, melindungi, dan memenuhi.
\end{itemize}

\subsection{Membangun Kesepakatan Bersama}

\subsubsection{Prinsip Dasar dan Kewarganegaraan}
\begin{itemize}
\item Musyawarah dilakukan untuk menyatukan berbagai pendapat guna mencapai kesepakatan bersama yang harus diterima dengan lapang dada dan dilaksanakan dengan penuh tanggung jawab.
\item Menurut Moh. Mahfud MD, kelemahan demokrasi yang mengedepankan prinsip menang-kalah berdasarkan suara terbanyak harus diimbangi dengan nomokrasi (negara hukum) agar tercipta keadilan dan kebenaran.
\item Menurut Mohammad Hatta, demokrasi dapat tertindas sementara karena kesalahannya sendiri, namun akan muncul kembali dengan keinsyafan setelah melalui cobaan.
\end{itemize}

\subsubsection{Keterbukaan Informasi Publik}
\begin{itemize}
\item Keterbukaan informasi publik merupakan kunci keterbukaan demokrasi dan kewajiban badan publik untuk mendorong tata kelola pemerintahan yang baik.
\item Masyarakat berhak mengakses informasi untuk berpartisipasi dalam pengambilan keputusan dan melakukan pengawasan terhadap kinerja pemerintah.
\item Informasi publik wajib dibuka secara transparan, kecuali informasi yang dikecualikan seperti rahasia negara, rahasia pribadi, dan rahasia bisnis.
\end{itemize}

\subsubsection{Penerapan Kesepakatan Bersama di Berbagai Lingkungan}
\begin{itemize}
\item Lingkungan Keluarga: Musyawarah digunakan untuk menyatukan perbedaan pendapat antaranggota keluarga, seperti menentukan kegiatan liburan atau keputusan penting keluarga lainnya.
\item Lingkungan Sekolah: Kesepakatan bersama dilakukan antara pimpinan sekolah, Komite Sekolah, guru, karyawan, dan siswa dalam menyusun tata tertib, pemilihan pengurus kelas atau OSIS, pembentukan kelompok piket, maupun agenda kegiatan sekolah.
\item Lingkungan Masyarakat: Kesepakatan bersama memperkuat integrasi sosial di tengah keberagaman suku, agama, dan budaya melalui pembentukan aturan RT atau RW maupun norma desa. Musyawarah Desa menjadi forum strategis dalam perencanaan pembangunan desa.
\item Lingkungan Bangsa dan Negara: Kesepakatan dilaksanakan oleh lembaga negara seperti MPR, DPR, DPD, dan DPRD dalam merumuskan peraturan perundang-undangan. Contoh historis utama adalah penetapan Pancasila dan UUD NRI Tahun 1945 sebagai dasar negara.
\end{itemize}
